\documentclass[11pt]{article}

\usepackage[utf8]{inputenc}
\usepackage[T1]{fontenc}
\usepackage{amsmath,amssymb,amsthm,mathtools}
\usepackage[margin=1.1in]{geometry}
\usepackage{booktabs}
\usepackage[colorlinks=true,linkcolor=blue,citecolor=blue]{hyperref}

\newtheorem{theorem}{Theorem}[section]
\newtheorem{proposition}[theorem]{Proposition}
\newtheorem{lemma}[theorem]{Lemma}
\newtheorem{corollary}[theorem]{Corollary}
\theoremstyle{definition}
\newtheorem{definition}[theorem]{Definition}
\newtheorem{example}[theorem]{Example}
\theoremstyle{remark}
\newtheorem{remark}[theorem]{Remark}
\newtheorem*{principleC}{Principle C (Contextual conservation)}

\newcommand{\Hi}{\mathcal{H}}
\newcommand{\PH}{P(\Hi)}
\newcommand{\F}{\mathcal{F}}
\newcommand{\C}{\mathcal{C}}
\newcommand{\Linf}{L_{\infty}}
\newcommand{\Om}{\Omega}
\newcommand{\Tr}{\operatorname{Tr}}

\title{Partition Logics, Their Continuum Extensions,\\
and the Conservation Law They Lack:\\
\large Four Clauses of Contextual Conservation}
\author{[Author]}
\date{June 2026}

\begin{document}
\maketitle

\begin{abstract}
Partition logics arise from pastings of Boolean algebras associated with
partitions of a set. We show that any continuum extension obtained by
pasting Boolean sub-$\sigma$-algebras of a common measurable space
$(\Om,\F)$ remains a \emph{concrete} logic: point evaluations furnish a
full, order-determining family of two-valued states, so by the
Kochen--Specker theorem the construction cannot be isomorphic to the
projection lattice $\PH$ for $\dim\Hi\ge 3$. Read positively, the no-go
isolates what such extensions lack: a \emph{conservation law for
probability across intertwined contexts}. We decompose this law into
four independent clauses---(C1) non-contextuality of value, (C2)
conservation in every context, (C3) continuous transitive symmetry on
pure states, (C4) rigidity of the conserved transition structure---each
enforced by a classical theorem (Kochen--Specker and Gleason for
(C1)--(C2) at dimension $\geq 3$; Hardy's continuity axiom for (C3);
Wigner's theorem and the Koecher--Vinberg classification with local
tomography for (C4)), and shown to be logically independent by a
scorecard of known theories, each obtained by deleting exactly one
clause. The dimension-two caveat is sharpened accordingly: the qubit
projection lattice is a horizontal sum, hence itself concrete, so at
$d=2$ the logical clauses do not bind and the entire quantum content of
the Bloch sphere---including the complementarity equality
$D^{2}+V^{2}=1$ and the Brukner--Zeilinger invariant---is carried by the
geometric clauses alone. A continuum limit of partition logics over a
classical point space therefore cannot approach quantum structure: the
continuation of partition logics is a change of category, not a limit.
\end{abstract}

\section{Introduction}

Partition logics arise in the study of empirical propositional structures,
finite automata, generalized urn models, and other contextual systems
\cite{Svozil1998,FoulisRandall1972,Wright1990,Svozil2005}: pastings of
Boolean algebras generated by partitions of a finite set, in which
propositions within one context are jointly decidable while propositions
from different contexts may be incompatible. A natural question is
whether a \emph{continuum} analogue---a pasting over a continuum state
space---could approach or reproduce the propositional structure of
quantum mechanics.

We answer this in two movements. \emph{Negatively}
(Section~\ref{sec:extension}): the continuum extension via Boolean
sub-$\sigma$-algebras of a common measurable space $(\Om,\F)$ is
\emph{concrete by construction}---point evaluations supply a full,
order-determining family of two-valued states, instantiating the
classical characterization of set-representable logics
\cite{Gudder1988,PtakPulmannova1991}---so by the Kochen--Specker theorem
\cite{KochenSpecker1967} it cannot be isomorphic to $\PH$ for
$\dim\Hi\ge 3$, with Gleason's theorem \cite{Gleason1957} an independent
route in the $\sigma$-additive setting. The obstruction is ontological:
retaining a common point space is retaining a hidden-variable model.

\emph{Positively} (Sections~\ref{sec:dimtwo}--\ref{sec:principle}): the
no-go is diagnostic. What the construction lacks is a \emph{conservation
law for probability across intertwined contexts}, which we decompose
into four independent clauses:
\begin{itemize}
\item[\textbf{(C1)}] \emph{Non-contextuality of value}: a proposition
shared by several contexts receives one value;
\item[\textbf{(C2)}] \emph{Conservation in every context}: the values
across each context sum to unity;
\item[\textbf{(C3)}] \emph{Continuous transitive symmetry}: pure states
form a connected manifold on which the reversible transformations act
continuously and transitively;
\item[\textbf{(C4)}] \emph{Budget rigidity}: the transition structure
between pure states is itself the conserved quantity, and the symmetries
are exactly those preserving it.
\end{itemize}
Each clause is enforced by a classical theorem---Kochen--Specker and
Gleason for (C1)--(C2) at dimension $\geq 3$, Hardy's continuity axiom
for (C3), Wigner and Koecher--Vinberg for (C4)---and the clauses are
logically independent (Section~\ref{subsec:scorecard}). Two observations
frame the analysis. First, the qubit lattice $P(\mathbb{C}^{2})$ is a
horizontal sum, hence concrete, and Bell's 1966 model \cite{Bell1966}
reproduces its probabilities; at $d=2$ the logical clauses do not bind,
and the quantum content of the Bloch sphere---the equality $D^2+V^2=1$
\cite{GreenbergerYasin1988,JaegerShimonyVaidman1995,Englert1996} and the
Brukner--Zeilinger invariant \cite{BruknerZeilinger1999,Zeilinger1999}---is
carried by the geometric clauses alone
(Section~\ref{subsec:qubitequation}). Second, cardinality is not the
missing ingredient: enlarging $\Om$ enlarges a simplex but never makes
it round, and concrete logics violate (C3) categorically
(Lemma~\ref{lem:discrete}).

\section{The Continuum Extension and Its Concreteness}\label{sec:extension}

Let $\Om$ be a set equipped with a $\sigma$-algebra $\F$. We work with
subsets of $\Om$ directly (not modulo null sets), so that pointwise
truth assignments are well defined.

\begin{definition}
A \emph{context} on $(\Om,\F)$ is a Boolean sub-$\sigma$-algebra
$\C\subseteq\F$. Given a family $\{\C_{\alpha}\}_{\alpha\in A}$ of
contexts, the associated \emph{concrete logic} is
$\Linf=\bigcup_{\alpha\in A}\C_{\alpha}$, partially ordered by
inclusion, with $A^{\perp}=\Om\setminus A$.
\end{definition}

In the finite case each context is the algebra of unions of blocks of a
partition, so this directly generalizes partition logic. We make no
claim that $\Linf$ is orthomodular---unions of Boolean
$\sigma$-algebras need not satisfy the orthomodular law, and ensuring it
requires delicate pasting conditions \cite{NavaraPtak1991}; our analysis
needs only that $\Linf$ is a family of subsets of a common point space.

\begin{example}\label{ex:sphere}
Let $\Om=S^{2}$ with its Borel $\sigma$-algebra. For each unit vector
$\mathbf{n}$, let $\C_{\mathbf n}$ be generated by the partition of the
sphere into open northern hemisphere, open southern hemisphere, and
equator. Since $\C_{-\mathbf n}=\C_{\mathbf n}$, the contexts are
indexed by $\mathbb{RP}^{2}$, and for $\mathbf m\neq\pm\mathbf n$ one
has $\C_{\mathbf n}\cap\C_{\mathbf m}=\{\emptyset,\Om\}$: a
\emph{horizontal sum}. We flag a coincidence for later use: the maximal
contexts of the \emph{qubit}, $\{P_{\mathbf n},P_{-\mathbf n}\}$, are
indexed by the same $\mathbb{RP}^{2}$. Example~\ref{ex:sphere} is the
\emph{classicalization} of the qubit---same sphere, same context index,
but with $S^{2}$ as \emph{sample space} rather than \emph{state space}.
\end{example}

For each $\omega\in\Om$ define the \emph{point evaluation}
$v_{\omega}(A)=1$ if $\omega\in A$, else $0$. Each
$v_{\omega}\colon\Linf\to\{0,1\}$ restricts to a Boolean homomorphism on
every context: a \emph{dispersion-free state}, sharp on every
proposition in every context simultaneously.

\begin{definition}
A family $\mathcal{S}$ of two-valued states on a poset $L$ is
\emph{order-determining} if
$\bigl(\forall s\in\mathcal S:\ s(A)\le s(B)\bigr)\implies A\le B$.
\end{definition}

\begin{proposition}\label{prop:od}
$\{v_{\omega}\}_{\omega\in\Om}$ is order-determining for $\Linf$.
\end{proposition}

\begin{proof}
If $v_\omega(A)\le v_\omega(B)$ for all $\omega$ but $A\not\subseteq B$,
pick $\omega_{0}\in A\setminus B$; then
$v_{\omega_{0}}(A)=1>0=v_{\omega_{0}}(B)$.
\end{proof}

This instantiates the classical characterization of concrete logics
\cite{Gudder1988,PtakPulmannova1991}: a logic is set-representable iff it
admits a full, order-determining set of two-valued states.

\begin{corollary}\label{cor:ks}
If $\dim\Hi\ge 3$, then $\Linf$ is not isomorphic to $\PH$.
\end{corollary}

\begin{proof}
Kochen--Specker \cite{KochenSpecker1967} excludes any two-valued state
on $\PH$ for $\dim\Hi\ge3$; $\Linf$ admits the $v_\omega$.
\end{proof}

\begin{remark}[Gleason]\label{rem:gleason}
Independently, Gleason's theorem \cite{Gleason1957} shows that for
$\dim\Hi\ge3$ every countably additive probability measure on $\PH$ is
$\Tr(\rho\,\cdot)$; in particular no $\{0,1\}$-valued countably additive
measure exists. The Kochen--Specker route is sharper, requiring no
additivity beyond finite additivity within contexts.
\end{remark}

The obstruction is thus ontological, not measure-theoretic: the points
of $\Om$ are hidden variables, and point evaluation is the deterministic
valuation.

\section{Dimension Two, and the Diagnosis}\label{sec:dimtwo}

\subsection{Where the obstruction does not apply}

\begin{proposition}\label{prop:qubit}
As an orthocomplemented poset, $P(\mathbb{C}^{2})$ is the horizontal sum
of continuum-many four-element Boolean algebras
($\mathrm{MO}_{\mathfrak c}$ in the notation of \cite{Kalmbach1983}\,)
and admits a full, order-determining family of two-valued states; hence
it is itself a concrete logic.
\end{proposition}

\begin{proof}[Proof sketch]
The nontrivial projections are rank one; two distinct ones commute iff
complementary. The blocks $\{0,P,P^{\perp},I\}$ intersect pairwise in
$\{0,I\}$ and are indexed by $\mathbb{RP}^{2}$. A two-valued state may
select, independently in each block, which of $P,P^{\perp}$ receives
$1$; additivity binds only within blocks, so all selections are
admissible, and the family is order-determining.
\end{proof}

This is the order-theoretic face of Bell's observation \cite{Bell1966}
that the qubit admits a hidden-variable model reproducing all quantum
probabilities: at dimension two the Kochen--Specker divide has not
opened. Note that Example~\ref{ex:sphere} and
Proposition~\ref{prop:qubit} exhibit horizontal sums over the
\emph{same} index set $\mathbb{RP}^{2}$: at $d=2$ the classicalization
of the qubit loses nothing \emph{logically}; what it loses is
metric---the transition probabilities between non-orthogonal states
(Section~\ref{subsec:qubitequation}). The structural moral: the
obstruction of Corollary~\ref{cor:ks} is driven not by the multiplicity
or cardinality of contexts but by their \emph{intertwining}---the
sharing of nontrivial propositions between maximal contexts, absent in
any horizontal sum and dense in $\PH$ for $\dim\Hi\ge3$.

\subsection{Cardinality is not the missing ingredient}

The probability measures on $(\Om,\F)$ restrict to states on $\Linf$,
forming a Choquet simplex whose extreme points are the $v_{\omega}$;
the reversible transformations compatible with the point structure are
bijections of $\Om$, which merely \emph{permute} the extreme points.

\begin{lemma}\label{lem:discrete}
For any path $t\mapsto\omega(t)$ in $\Om$ and $A\in\Linf$, the function
$t\mapsto v_{\omega(t)}(A)\in\{0,1\}$ is constant or discontinuous; and
every $v_{\omega}$ is maximally sharp in every context simultaneously.
\end{lemma}

\begin{proof}
Immediate: the function is $\{0,1\}$-valued, and $v_\omega$ restricts to
a homomorphism on each $\C_\alpha$.
\end{proof}

Trivial as it is, Lemma~\ref{lem:discrete} states precisely what a
classical pasting can never do: realize a \emph{continuous trade-off} in
which sharpening the answers in one context continuously degrades those
in another. Dispersion-free states have an unbounded information
budget---they saturate all contexts at once---and the only reversible
dynamics is permutation. Enlarging $\Om$ from a finite set to a
continuum enlarges the simplex; it does not bend it.

Example~\ref{ex:sphere} now appears in its true light:

\begin{center}
\begin{tabular}{@{}lll@{}}
\toprule
 & \textbf{Example~\ref{ex:sphere}} & \textbf{Qubit (Bloch sphere)}\\
\midrule
Role of $S^{2}$ & sample space & pure-state space\\
A point is & a complete answer-book & a budget of potential answers\\
Rotations act on & points and contexts & states, conserving the budget\\
Sharpness & maximal in all contexts & maximal in one context only\\
Conserved quantity & none (vacuous) & $S_{1}^{2}+S_{2}^{2}+S_{3}^{2}=1$\\
\bottomrule
\end{tabular}
\end{center}

The right-hand conserved quantity is the Brukner--Zeilinger operational
information invariant \cite{BruknerZeilinger1999}, which for
path/interference observables specializes to $D^{2}+V^{2}=1$
\cite{GreenbergerYasin1988,JaegerShimonyVaidman1995,Englert1996}. But by
Proposition~\ref{prop:qubit}, at dimension two this conservation law
does not yet exclude dispersion-free valuations: it binds mixed states
and continuous symmetries, not blockwise two-valued homomorphisms. The
law acquires modal force---excluding things---only when contexts
intertwine. The next section makes its anatomy explicit.

\section{Principle C: Four Clauses of Contextual Conservation}\label{sec:principle}

\begin{definition}\label{def:budget}
Let $L$ be a pasting of Boolean ($\sigma$-)algebras $\{\C_\alpha\}$, not
necessarily concrete. A \emph{probability budget} on $L$ is a map
$\mu\colon L\to[0,1]$ with $\mu(1)=1$, (finitely or countably) additive
within each context.
\end{definition}

Noncontextuality is built in: $\mu$ is defined on propositions, not
(proposition, context) pairs. On a horizontal sum the consistency
requirement is vacuous; on a concrete logic it is satisfiable in
$\{0,1\}$---by Proposition~\ref{prop:od}, lavishly so. The entire
content of the quantum/classical divide is the behaviour of budgets
under \emph{intertwining}.

\begin{principleC}
A \emph{state} of a contextual theory is a probability budget subject to:
\begin{itemize}
\item[\textbf{(C1)}] \emph{Non-contextuality of value.} For every
proposition $a\in L$ and all contexts $\C,\C'$ containing $a$, the value
$\mu(a)$ is the same whether computed in $\C$ or in $\C'$.
\item[\textbf{(C2)}] \emph{Conservation in every context.} For every
context $\C$ with atomic decomposition $\{a_{1},a_{2},\dots\}$,
$\sum_{i}\mu(a_{i})=1$, with (countable) additivity within $\C$.
\item[\textbf{(C3)}] \emph{Continuous transitive symmetry.} The pure
states form a connected topological manifold $\mathcal{P}$, and the
group $G$ of reversible transformations acts continuously and
transitively on $\mathcal{P}$.
\item[\textbf{(C4)}] \emph{Budget rigidity.} There is a continuous
transition probability $p\colon\mathcal{P}\times\mathcal{P}\to[0,1]$,
with $p(a,b)=1$ iff $a=b$ and $p(a,b)=0$ iff $a\perp b$, such that $G$
is \emph{exactly} the group of bijections of $\mathcal{P}$ preserving
$p$; equivalently, the conserved invariant (purity; for the qubit, the
squared Bloch length) determines the symmetry group, and conversely.
\end{itemize}
\end{principleC}

Clauses (C1)--(C2) are \emph{logical} (the budget as a function on the
pasted contexts); clauses (C3)--(C4) are \emph{geometric} (the space of
budgets and its symmetries). The dimension-two caveat is precisely the
statement that the logical clauses can be satisfied classically at
$d=2$ while the geometric clauses already cannot.

\subsection{Clause-wise enforcement}\label{subsec:rigidity}

On $\PH$ with $\dim\Hi\ge3$ the rigidity of Principle C decomposes
clause-wise:

\begin{enumerate}
\item \textbf{(C1)$+$(C2), two-valued: impossible} (Kochen--Specker
\cite{KochenSpecker1967}). A $\{0,1\}$-budget would have zero entropy in
every context simultaneously; dense intertwining makes this
combinatorially impossible. Contrapositively, a concrete logic is
exactly one whose (C1)--(C2) constraints are slack enough to admit
saturation everywhere (Section~\ref{sec:extension}).
\item \textbf{(C1)$+$(C2), $[0,1]$-valued: unique} (Gleason
\cite{Gleason1957}). Consistency across intertwined contexts alone
forces $\mu=\Tr(\rho\,\cdot)$: the Born rule is the unique solution of
the logical clauses in the $\sigma$-additive setting, with no assumption
about dynamics, symmetry, or continuity.
\item \textbf{(C3): the classical/quantum separator} (Hardy
\cite{Hardy2001}). In Hardy's reconstruction, classical probability and
quantum theory differ in exactly one axiom: the existence of continuous
reversible transformations between pure states---clause (C3).
Lemma~\ref{lem:discrete} shows no concrete logic can satisfy it: over a
point space, pure states are vertices and reversible maps are
permutations.
\item \textbf{(C4): the selector of complex quantum theory} (Wigner
\cite{Bargmann1964}; Koecher--Vinberg \cite{FarautKoranyi1994}).
Wigner's theorem: the bijections of pure states preserving transition
probabilities are exactly the (anti)unitaries, as (C4) demands.
Homogeneity of the state cone (from (C3)) together with self-duality
(the (C4) pairing of states with effects) characterizes the formally
real Jordan algebras; adjoining local tomography selects complex
quantum theory
\cite{Hardy2001,DakicBrukner2011,MasanesMuller2011,%
ChiribellaDArianoPerinotti2011}. That the complex case is empirically
distinguishable has recently been demonstrated: real-Hilbert-space
quantum theory is falsified in network Bell scenarios \cite{Renou2021}.
\end{enumerate}

\subsection{Independence: a scorecard}\label{subsec:scorecard}

The clauses are logically independent; the proof is a museum of known
theories, each deleting exactly one clause.

\begin{center}
\small
\begin{tabular}{@{}lcccc@{}}
\toprule
\textbf{Theory} & \textbf{(C1)} & \textbf{(C2)} & \textbf{(C3)} &
\textbf{(C4)}\\
\midrule
Classical probability; concrete logics over $\Om$ & \checkmark &
\checkmark & --- & ---\\
Spekkens' toy theory \cite{Spekkens2007} & \checkmark & \checkmark &
(discrete only) & ---\\
Real-Hilbert-space quantum theory & \checkmark & \checkmark &
\checkmark & ---\\
Jordan-algebraic state spaces (e.g.\ exceptional) & \checkmark &
\checkmark & \checkmark & ---\\
Complex quantum theory, $\dim\Hi\geq 3$ & \checkmark & \checkmark &
\checkmark & \checkmark\\
\bottomrule
\end{tabular}
\end{center}

\begin{remark}[What each deletion costs]\label{rem:deletions}
Deleting (C3) returns the simplex: pure states discrete, symmetries
permutations, no interference---the entire family of partition logics
and their continuum extensions (Lemma~\ref{lem:discrete}). Keeping only
a discrete shadow of (C3) yields Spekkens' toy theory
\cite{Spekkens2007}: an epistemically restricted classical theory
satisfying (C1)--(C2) that reproduces complementarity, no-cloning, and
teleportation analogues---calibrating how much apparent quantumness
resides in the logical clauses alone---but with no Bell violation, no
continuous interference phase, no Gleason rigidity. Deleting (C4) while
keeping (C3) yields the real and Jordan-algebraic alternatives,
empirically distinguishable from the complex case and now excluded
\cite{Renou2021}. Deleting (C1) admits contextual value assignments and
exits the framework; deleting (C2) severs the connection to operational
probability. The conjunction carves out complex quantum theory, and no
clause is redundant.
\end{remark}

\subsection{The qubit equation: four clauses in one formula}\label{subsec:qubitequation}

For the qubit, Principle C condenses into a single equation. For pure
states $\psi,\phi$ at Bloch angle $\theta$,
\begin{equation}\label{eq:malus}
p(\psi,\phi)\;=\;\bigl|\langle\psi|\phi\rangle\bigr|^{2}
\;=\;\cos^{2}\!\frac{\theta}{2}.
\end{equation}
Each clause is legible in \eqref{eq:malus}: \textbf{(C1)} the value
depends on the pair of states only; \textbf{(C2)}
$\cos^{2}(\theta/2)+\sin^{2}(\theta/2)=1$ in any basis
$\{\phi,\phi^{\perp}\}$; \textbf{(C3)} $\theta$ varies continuously and
$SO(3)$ acts transitively; \textbf{(C4)} $\cos^{2}(\theta/2)$ is, up to
trivial reparametrization, the unique transition probability on $S^{2}$
admitting a transitive continuous invariance group of the Wigner type
(cf.\ \cite{BruknerZeilingerMalus}).

Two consequences follow. First, $D^{2}+V^{2}=1$ and the
Brukner--Zeilinger invariant are \eqref{eq:malus} read along orthogonal
great circles: the (C3)--(C4) content of the qubit. Second,
Proposition~\ref{prop:qubit} shows the (C1)--(C2) content of the qubit
to be classically satisfiable, and Example~\ref{ex:sphere} shows its
\emph{logic} can be carried by a sample space; what no sample space can
carry is \eqref{eq:malus} itself, since by Lemma~\ref{lem:discrete}
point evaluations support only the Kronecker structure
$p(\omega,\omega')=\delta_{\omega\omega'}$. The quantum content of the
Bloch sphere is not its logic but its \emph{metric}. Textbook
discussions of complementarity conducted entirely at $d=2$ probe the
conservation law only through its geometric clauses---which is precisely
why it has so often been mistaken there for a triviality of
normalization.

\section{Conclusion: A Category Change, Not a Limit}

The continuum extension of partition logics is faithful to the finite
theory in every respect but one: it inherits, and cannot shed, the
underlying point space. This single inheritance is fatal to any
convergence toward quantum structure (Proposition~\ref{prop:od},
Corollary~\ref{cor:ks}); the failure is invisible at dimension two
(Proposition~\ref{prop:qubit}); and the correct diagnosis is the absence
of a conservation law for probability across intertwined contexts,
resolved here into four independent clauses. The logical clauses
(C1)--(C2) are enforced at dimension $\geq3$ by Kochen--Specker and
Gleason; the geometric clauses (C3)--(C4) by Hardy's continuity axiom
and by Wigner's theorem with the Koecher--Vinberg classification and
local tomography. Deleting (C3) yields classical probability and every
partition logic; weakening it to a discrete action yields Spekkens'
toy theory; deleting (C4) yields the real and Jordan-algebraic
alternatives, now empirically excluded \cite{Renou2021}. Their
conjunction characterizes complex quantum theory. At dimension two the
logical clauses are slack and the entire divide is carried by the
metric \eqref{eq:malus}, not the logic---which is why Bloch-sphere
complementarity could be mistaken for a triviality of normalization.

A continuum limit of partition logics built on a single classical point
space will always admit a complete answer-book behind the contexts, and
a theory with an answer-book has nothing to conserve: its budgets
saturate every context at once, its symmetries permute, and its
transition structure is the Kronecker delta. The continuation of
partition logics is therefore not a limit within the classical category
but an exchange of its basic object: trade the hidden answer-book for
the conserved budget. What remains of Bohr's celebrated analogy between
complementarity and relativity is then exactly this---not the relativity
of descriptions over a fixed stage, but a four-clause conservation law
that holds precisely because no stage is left for it to play on.

\begin{thebibliography}{99}

\bibitem{Svozil1998}
K.~Svozil, \emph{Quantum Logic}, Springer, Singapore, 1998.

\bibitem{FoulisRandall1972}
D.~J. Foulis and C.~H. Randall, ``Operational statistics. I. Basic
concepts,'' \emph{J. Math. Phys.} \textbf{13} (1972), 1667--1675.

\bibitem{Wright1990}
R.~Wright, ``Generalized urn models,'' \emph{Found. Phys.} \textbf{20}
(1990), 881--903.

\bibitem{Svozil2005}
K.~Svozil, ``Logical equivalence between generalized urn models and
finite automata,'' \emph{Int. J. Theor. Phys.} \textbf{44} (2005),
745--754.

\bibitem{Gudder1988}
S.~Gudder, \emph{Quantum Probability}, Academic Press, Boston, 1988.

\bibitem{PtakPulmannova1991}
P.~Pt\'ak and S.~Pulmannov\'a, \emph{Orthomodular Structures as Quantum
Logics}, Kluwer, Dordrecht, 1991.

\bibitem{NavaraPtak1991}
M.~Navara and P.~Pt\'ak, ``Almost Boolean orthomodular posets,''
\emph{J. Pure Appl. Algebra} \textbf{60} (1991), 105--111.

\bibitem{Kalmbach1983}
G.~Kalmbach, \emph{Orthomodular Lattices}, Academic Press, London, 1983.

\bibitem{KochenSpecker1967}
S.~Kochen and E.~P. Specker, ``The problem of hidden variables in quantum
mechanics,'' \emph{J. Math. Mech.} \textbf{17} (1967), 59--87.

\bibitem{Gleason1957}
A.~M. Gleason, ``Measures on the closed subspaces of a Hilbert space,''
\emph{J. Math. Mech.} \textbf{6} (1957), 885--893.

\bibitem{Bell1966}
J.~S. Bell, ``On the problem of hidden variables in quantum mechanics,''
\emph{Rev. Mod. Phys.} \textbf{38} (1966), 447--452.

\bibitem{Bargmann1964}
V.~Bargmann, ``Note on Wigner's theorem on symmetry operations,''
\emph{J. Math. Phys.} \textbf{5} (1964), 862--868.

\bibitem{GreenbergerYasin1988}
D.~M. Greenberger and A.~Yasin, ``Simultaneous wave and particle
knowledge in a neutron interferometer,'' \emph{Phys. Lett. A}
\textbf{128} (1988), 391--394.

\bibitem{JaegerShimonyVaidman1995}
G.~Jaeger, A.~Shimony, and L.~Vaidman, ``Two interferometric
complementarities,'' \emph{Phys. Rev. A} \textbf{51} (1995), 54--67.

\bibitem{Englert1996}
B.-G. Englert, ``Fringe visibility and which-way information: an
inequality,'' \emph{Phys. Rev. Lett.} \textbf{77} (1996), 2154--2157.

\bibitem{BruknerZeilinger1999}
\v{C}.~Brukner and A.~Zeilinger, ``Operationally invariant information in
quantum measurements,'' \emph{Phys. Rev. Lett.} \textbf{83} (1999),
3354--3357.

\bibitem{Zeilinger1999}
A.~Zeilinger, ``A foundational principle for quantum mechanics,''
\emph{Found. Phys.} \textbf{29} (1999), 631--643.

\bibitem{BruknerZeilingerMalus}
\v{C}.~Brukner and A.~Zeilinger, ``Malus' law and quantum information,''
\emph{Acta Phys. Slovaca} \textbf{49} (1999), 647--652.

\bibitem{Spekkens2007}
R.~W. Spekkens, ``Evidence for the epistemic view of quantum states: a
toy theory,'' \emph{Phys. Rev. A} \textbf{75} (2007), 032110.

\bibitem{Hardy2001}
L.~Hardy, ``Quantum theory from five reasonable axioms,''
arXiv:quant-ph/0101012 (2001).

\bibitem{DakicBrukner2011}
B.~Daki\'c and \v{C}.~Brukner, ``Quantum theory and beyond: is
entanglement special?'' in \emph{Deep Beauty}, H.~Halvorson (ed.),
Cambridge University Press, 2011, 365--392.

\bibitem{MasanesMuller2011}
L.~Masanes and M.~P. M\"uller, ``A derivation of quantum theory from
physical requirements,'' \emph{New J. Phys.} \textbf{13} (2011), 063001.

\bibitem{ChiribellaDArianoPerinotti2011}
G.~Chiribella, G.~M. D'Ariano, and P.~Perinotti, ``Informational
derivation of quantum theory,'' \emph{Phys. Rev. A} \textbf{84} (2011),
012311.

\bibitem{FarautKoranyi1994}
J.~Faraut and A.~Kor\'anyi, \emph{Analysis on Symmetric Cones}, Oxford
University Press, Oxford, 1994.

\bibitem{Renou2021}
M.-O. Renou, D.~Trillo, M.~Weilenmann, T.~P. Le, A.~Tavakoli, N.~Gisin,
A.~Ac\'in, and M.~Navascu\'es, ``Quantum theory based on real numbers can
be experimentally falsified,'' \emph{Nature} \textbf{600} (2021),
625--629.

\end{thebibliography}

\end{document}
